\subsection{Packed polynomial and the tensor-valued evaluation}\label{sec:dpr:eval}
Fix a Boolean witness $w:I_\ell\times I_h\to\F_2$,
with $h\geq1$, and the evaluation points
$r_{\rm lo}\in L^\ell$, $r_{\rm hi}\in L^h$
from Chapter~7. Write $q_v=\operatorname{eq}_\ell
(v,r_{\rm lo})$, $h_j=\operatorname{eq}_h(j,r_{\rm hi})$.

\begin{definition}[Packed multilinear polynomial]\label{def:dpr:packed-poly}
The packed table is
$t'(j):=\sum_v\beta_vw_{v,j}=\widehat w_j$,
$j\in I_h$. Its multilinear extension
is
\[
 T'(X)=\sum_{j\in I_h}\widehat w_j
 \operatorname{eq}_h(j,X)\in
 L[X_0,\ldots,X_{h-1}]_{\mathrm{mult}}.
\]
Let $\varphi_1(T')$ denote the polynomial
obtained by applying $\varphi_1$ to each
coefficient of $T'$.
\end{definition}

\begin{lemma}[Tensor expansion on Boolean columns]\label{lem:dpr:tensor-column}
For every $j\in I_h$,
\[
 \varphi_1(\widehat w_j)
 =\sum_v\varphi_0(w_{v,j})\varphi_1(\beta_v).
\]
\end{lemma}
\begin{proof}
By definition
$\widehat w_j=\sum_v\beta_vw_{v,j}$.
Apply the additive ring homomorphism
$\varphi_1$.
Because $w_{v,j}\in\F_2$, the scalar
embeddings agree:
$\varphi_1(w_{v,j})=\varphi_0(w_{v,j})$.
Rearranging the commutative products
gives the displayed expansion.
\end{proof}

\begin{definition}[Tensor partial evaluation]\label{def:dpr:tensor-S}
Put
\[
 S:=\varphi_1(T')(
 \varphi_0(r_{{\rm hi},0}),\ldots,
 \varphi_0(r_{{\rm hi},h-1}))\in A.
\]
\end{definition}

\begin{lemma}[Evaluation of equality weights through the left embedding]\label{lem:dpr:eq-embedded}
For each $j\in I_h$,
\[
 \operatorname{eq}_h(j,\varphi_0(r_{\rm hi}))
 =\varphi_0(h_j).
\]
\end{lemma}
\begin{proof}
Each equality weight is a product of
factors of the form $(1-r_i)$ or $r_i$,
selected by the Boolean index $j_i$.
A unital ring homomorphism preserves
addition, products, and the elements
$0,1$. Applying $\varphi_0$
coordinatewise to the defining product
thus gives the same result as evaluating
at the embedded point.
\end{proof}

\begin{theorem}[Columns of the tensor evaluation]\label{thm:dpr:column-evaluation}
For every $v\in I_\ell$,
\[
 \boxed{
 \operatorname{Col}(S)_v
  =\sum_{j\in I_h}h_jw_{v,j}
  =s_v^*.}
\]
\end{theorem}
\begin{proof}
Multilinear interpolation and
Lemma~\ref{lem:dpr:eq-embedded} give
\[
 S=\sum_j\varphi_1(\widehat w_j)
                 \varphi_0(h_j).
\]
Insert Lemma~\ref{lem:dpr:tensor-column}:
\[
 S=\sum_{j,v}\varphi_0(w_{v,j}h_j)
                       \varphi_1(\beta_v).
\]
Interchange the finite sums, use
additivity of $\varphi_0$, and obtain
\[
 S=\sum_v
     \varphi_0\left(\sum_jh_jw_{v,j}\right)
     \varphi_1(\beta_v).
\]
Uniqueness of column coefficients in
Lemma~\ref{lem:dpr:column-basis}
shows that the $v$-th coefficient
is $\sum_jh_jw_{v,j}$, precisely the
partial evaluation $s_v^*$ defined
in Chapter~7.
\end{proof}

\begin{theorem}[Original evaluation is a column contraction]\label{thm:dpr:claim-contraction}
The true evaluation of the original
Boolean multilinear table is
\[
 \boxed{y_w=C_q(S).}
\]
\end{theorem}
\begin{proof}
By Theorem~\ref{thm:dpr:column-evaluation},
$C_q(S)=\sum_vq_vs_v^*$.
Lemma~\ref{lem:pack:partial-eval}
identifies this finite sum with $y_w$.
\end{proof}

\begin{lemma}[Multiplication recovers only the packed high-axis evaluation]\label{lem:dpr:mu-evaluation}
\[
 \mu(S)=T'(r_{\rm hi})
       =\sum_v\beta_vs_v^*.
\]
\end{lemma}
\begin{proof}
Apply the multiplication map $\mu$
to the displayed expansion of $S$
in the proof of
Theorem~\ref{thm:dpr:column-evaluation}.
As $\mu(\varphi_0(a)\varphi_1(b))=ab$,
the image is $\sum_v\beta_vs_v^*$.
On the other hand multilinear
interpolation gives
$T'(r_{\rm hi})=\sum_jh_j\widehat w_j
=\sum_{j,v}h_j\beta_vw_{v,j}
=\sum_v\beta_vs_v^*$.
\end{proof}

\begin{definition}[Sumcheck summand over $A$]\label{def:dpr:H}
Define the $A$-coefficient polynomial
\[
 H(X):=\varphi_1(T')(X)\,
   \operatorname{eq}_h(\varphi_0(r_{\rm hi}),X)
 \in A[X_0,\ldots,X_{h-1}],
\]
where the multivariate equality polynomial
uses the usual affine factors in each variable.
\end{definition}

\begin{lemma}[Individual variable degree of the sumcheck summand]\label{lem:dpr:H-degree}
The degree of $H$ in each variable $X_i$
is at most $2$.
\end{lemma}
\begin{proof}
The polynomial $T'$ is multilinear,
so its coefficientwise image under
$\varphi_1$ has degree at most $1$
in each variable. The equality
polynomial is also multilinear, hence
has individual variable degree at most
$1$. In a commutative polynomial ring,
degrees of products in any fixed variable
are bounded above by the sum of the
degrees of the factors. Therefore
$\deg_{X_i}H\leq2$.
\end{proof}

\begin{theorem}[Tensor Boolean sum equals its partial evaluation]\label{thm:dpr:boolean-sum}
\[
 \boxed{
 \sum_{j\in I_h}H(j)=S.}
\]
\end{theorem}
\begin{proof}
For Boolean $j$, multilinear interpolation and
Lemma~\ref{lem:dpr:eq-embedded} give
\[
 \varphi_1(T')(j)=\varphi_1(\widehat w_j),
 \qquad
 \operatorname{eq}_h(\varphi_0(r_{\rm hi}),j)=\varphi_0(h_j).
\]
Consequently,
\[
 \sum_jH(j)
 =\sum_j\varphi_1(\widehat w_j)\varphi_0(h_j)=S,
\]
where the last equality is the expansion in
Theorem~\ref{thm:dpr:column-evaluation}.
\end{proof}

\begin{lemma}[Reduction to a large-field evaluation at a common challenge]\label{lem:dpr:terminal}
For $z\in L^h$,
\[
 \boxed{
 H(\varphi_1(z))
 =\varphi_1(T'(z))\,
 \operatorname{eq}_h(\varphi_0(r_{\rm hi}),\varphi_1(z)).
 }
\]
\end{lemma}
\begin{proof}
Evaluate Definition~\ref{def:dpr:H}
at the tuple of elements $\varphi_1(z_i)$.
Because $\varphi_1$ is a ring homomorphism,
substituting $\varphi_1(z_i)$ into the
coefficientwise image $\varphi_1(T')$
equals $\varphi_1(T'(z))$. The second
factor is already the displayed public
equality polynomial at that tuple.
\end{proof}

\begin{remark}[Why the tensor is essential]
The ordinary packed value $T'(r_{\rm hi})$
equals $\mu(S)$; it does not preserve
the independent columns of $S$.
The verifier needs these columns to
apply its arbitrary low-axis weights $q$.
The noninjectivity of $\mu$ from
Lemma~\ref{lem:dpr:multiplication-loss}
explains why evaluating only $T'$ at
$r_{\rm hi}$ cannot replace the tensor
phase of the original compiler.
\end{remark}
